% Options for packages loaded elsewhere
\PassOptionsToPackage{unicode}{hyperref}
\PassOptionsToPackage{hyphens}{url}
%
\documentclass[
  ignorenonframetext,
]{beamer}
\usepackage{pgfpages}
\setbeamertemplate{caption}[numbered]
\setbeamertemplate{caption label separator}{: }
\setbeamercolor{caption name}{fg=normal text.fg}
\beamertemplatenavigationsymbolsempty
% Prevent slide breaks in the middle of a paragraph
\widowpenalties 1 10000
\raggedbottom
\setbeamertemplate{part page}{
  \centering
  \begin{beamercolorbox}[sep=16pt,center]{part title}
    \usebeamerfont{part title}\insertpart\par
  \end{beamercolorbox}
}
\setbeamertemplate{section page}{
  \centering
  \begin{beamercolorbox}[sep=12pt,center]{part title}
    \usebeamerfont{section title}\insertsection\par
  \end{beamercolorbox}
}
\setbeamertemplate{subsection page}{
  \centering
  \begin{beamercolorbox}[sep=8pt,center]{part title}
    \usebeamerfont{subsection title}\insertsubsection\par
  \end{beamercolorbox}
}
\AtBeginPart{
  \frame{\partpage}
}
\AtBeginSection{
  \ifbibliography
  \else
    \frame{\sectionpage}
  \fi
}
\AtBeginSubsection{
  \frame{\subsectionpage}
}
\usepackage{amsmath,amssymb}
\usepackage{lmodern}
\usepackage{iftex}
\ifPDFTeX
  \usepackage[T1]{fontenc}
  \usepackage[utf8]{inputenc}
  \usepackage{textcomp} % provide euro and other symbols
\else % if luatex or xetex
  \usepackage{unicode-math}
  \defaultfontfeatures{Scale=MatchLowercase}
  \defaultfontfeatures[\rmfamily]{Ligatures=TeX,Scale=1}
\fi
\usetheme[]{Montpellier}
% Use upquote if available, for straight quotes in verbatim environments
\IfFileExists{upquote.sty}{\usepackage{upquote}}{}
\IfFileExists{microtype.sty}{% use microtype if available
  \usepackage[]{microtype}
  \UseMicrotypeSet[protrusion]{basicmath} % disable protrusion for tt fonts
}{}
\makeatletter
\@ifundefined{KOMAClassName}{% if non-KOMA class
  \IfFileExists{parskip.sty}{%
    \usepackage{parskip}
  }{% else
    \setlength{\parindent}{0pt}
    \setlength{\parskip}{6pt plus 2pt minus 1pt}}
}{% if KOMA class
  \KOMAoptions{parskip=half}}
\makeatother
\usepackage{xcolor}
\IfFileExists{xurl.sty}{\usepackage{xurl}}{} % add URL line breaks if available
\IfFileExists{bookmark.sty}{\usepackage{bookmark}}{\usepackage{hyperref}}
\hypersetup{
  hidelinks,
  pdfcreator={LaTeX via pandoc}}
\urlstyle{same} % disable monospaced font for URLs
\newif\ifbibliography
\usepackage{color}
\usepackage{fancyvrb}
\newcommand{\VerbBar}{|}
\newcommand{\VERB}{\Verb[commandchars=\\\{\}]}
\DefineVerbatimEnvironment{Highlighting}{Verbatim}{commandchars=\\\{\}}
% Add ',fontsize=\small' for more characters per line
\usepackage{framed}
\definecolor{shadecolor}{RGB}{248,248,248}
\newenvironment{Shaded}{\begin{snugshade}}{\end{snugshade}}
\newcommand{\AlertTok}[1]{\textcolor[rgb]{0.94,0.16,0.16}{#1}}
\newcommand{\AnnotationTok}[1]{\textcolor[rgb]{0.56,0.35,0.01}{\textbf{\textit{#1}}}}
\newcommand{\AttributeTok}[1]{\textcolor[rgb]{0.77,0.63,0.00}{#1}}
\newcommand{\BaseNTok}[1]{\textcolor[rgb]{0.00,0.00,0.81}{#1}}
\newcommand{\BuiltInTok}[1]{#1}
\newcommand{\CharTok}[1]{\textcolor[rgb]{0.31,0.60,0.02}{#1}}
\newcommand{\CommentTok}[1]{\textcolor[rgb]{0.56,0.35,0.01}{\textit{#1}}}
\newcommand{\CommentVarTok}[1]{\textcolor[rgb]{0.56,0.35,0.01}{\textbf{\textit{#1}}}}
\newcommand{\ConstantTok}[1]{\textcolor[rgb]{0.00,0.00,0.00}{#1}}
\newcommand{\ControlFlowTok}[1]{\textcolor[rgb]{0.13,0.29,0.53}{\textbf{#1}}}
\newcommand{\DataTypeTok}[1]{\textcolor[rgb]{0.13,0.29,0.53}{#1}}
\newcommand{\DecValTok}[1]{\textcolor[rgb]{0.00,0.00,0.81}{#1}}
\newcommand{\DocumentationTok}[1]{\textcolor[rgb]{0.56,0.35,0.01}{\textbf{\textit{#1}}}}
\newcommand{\ErrorTok}[1]{\textcolor[rgb]{0.64,0.00,0.00}{\textbf{#1}}}
\newcommand{\ExtensionTok}[1]{#1}
\newcommand{\FloatTok}[1]{\textcolor[rgb]{0.00,0.00,0.81}{#1}}
\newcommand{\FunctionTok}[1]{\textcolor[rgb]{0.00,0.00,0.00}{#1}}
\newcommand{\ImportTok}[1]{#1}
\newcommand{\InformationTok}[1]{\textcolor[rgb]{0.56,0.35,0.01}{\textbf{\textit{#1}}}}
\newcommand{\KeywordTok}[1]{\textcolor[rgb]{0.13,0.29,0.53}{\textbf{#1}}}
\newcommand{\NormalTok}[1]{#1}
\newcommand{\OperatorTok}[1]{\textcolor[rgb]{0.81,0.36,0.00}{\textbf{#1}}}
\newcommand{\OtherTok}[1]{\textcolor[rgb]{0.56,0.35,0.01}{#1}}
\newcommand{\PreprocessorTok}[1]{\textcolor[rgb]{0.56,0.35,0.01}{\textit{#1}}}
\newcommand{\RegionMarkerTok}[1]{#1}
\newcommand{\SpecialCharTok}[1]{\textcolor[rgb]{0.00,0.00,0.00}{#1}}
\newcommand{\SpecialStringTok}[1]{\textcolor[rgb]{0.31,0.60,0.02}{#1}}
\newcommand{\StringTok}[1]{\textcolor[rgb]{0.31,0.60,0.02}{#1}}
\newcommand{\VariableTok}[1]{\textcolor[rgb]{0.00,0.00,0.00}{#1}}
\newcommand{\VerbatimStringTok}[1]{\textcolor[rgb]{0.31,0.60,0.02}{#1}}
\newcommand{\WarningTok}[1]{\textcolor[rgb]{0.56,0.35,0.01}{\textbf{\textit{#1}}}}
\setlength{\emergencystretch}{3em} % prevent overfull lines
\providecommand{\tightlist}{%
  \setlength{\itemsep}{0pt}\setlength{\parskip}{0pt}}
\setcounter{secnumdepth}{-\maxdimen} % remove section numbering
%\documentclass[unicode,10pt]{beamer}% 'unicode'が必要
\usepackage{luatexja}% 日本語を使う場合は必須
%\usepackage[japanese]{babel}	
%\usefonttheme[onlymath]{serif}
%\usepackage[T1]{fontenc} %これを使うと、ギリシャ文字が出ない
%\usepackage{textcomp}
%\usepackage[scale=1.0]{tgheros} % Sans serif font
%\usepackage[scaled]{beramono} % Monospace font
%\usepackage{luatexja-otf}
%\renewcommand{\kanjifamilydefault}{\gtdefault}
%\usepackage[haranoaji]{luatexja-preset}

% スライド番号の表示 %%%%%%%%%%%%%%%%%%%
\makeatletter
\setbeamertemplate{footline}{%
  \leavevmode%
  \hbox{%
  \begin{beamercolorbox}[wd=.5\paperwidth,ht=2.25ex,dp=1ex,right]{date in head/foot}%
    \insertframenumber{} \hspace*{2ex} % new version without total frames
  \end{beamercolorbox}}%
  \vskip0pt%
}
\makeatother
% スライド番号の表示 %%%%%%%%%%%%%%%%%%%
\ifLuaTeX
  \usepackage{selnolig}  % disable illegal ligatures
\fi

\title{第3章: 測定\\
(3.1. 戦時における民間人の被害を測定する)}
\subtitle{今井耕介 著\\
『社会科学のためのデータ分析入門 (QSS)』}
\author{}
\date{\vspace{-2.5em}2026-03-09}

\begin{document}
\frame{\titlepage}

\hypertarget{ux6226ux6642ux306bux304aux3051ux308bux6c11ux9593ux4ebaux306eux88abux5bb3ux3092ux6e2cux5b9aux3059ux308b}{%
\section{3.1
戦時における民間人の被害を測定する}\label{ux6226ux6642ux306bux304aux3051ux308bux6c11ux9593ux4ebaux306eux88abux5bb3ux3092ux6e2cux5b9aux3059ux308b}}

\begin{frame}[fragile]{アフガニスタンでの世論調査}
\protect\hypertarget{ux30a2ux30d5ux30acux30cbux30b9ux30bfux30f3ux3067ux306eux4e16ux8ad6ux8abfux67fb}{}
\begin{itemize}
\tightlist
\item
  \textbf{背景}:
  2001年以来の紛争下にあるアフガニスタンで、民間人の被害実態を調査。
\item
  \textbf{目的}:

  \begin{itemize}
  \tightlist
  \item
    戦闘員（ISAF:
    国際治安支援部隊、および反政府勢力：タリバン）による直接的な被害を測定する。
  \item
    被害が人々の政治的態度にどう影響するかを分析する。
  \end{itemize}
\item
  \textbf{データ}: 2010〜2011年に実施された全国的な対面調査データ
  (\texttt{afghan.csv})。
\end{itemize}
\end{frame}

\hypertarget{ux30c7ux30fcux30bfux306eux8aadux307fux8fbcux307fux3068ux57faux672cux7d71ux8a08ux91cf}{%
\section{3.1.1
データの読み込みと基本統計量}\label{ux30c7ux30fcux30bfux306eux8aadux307fux8fbcux307fux3068ux57faux672cux7d71ux8a08ux91cf}}

\begin{frame}[fragile]{afghan データの読み込み (1)}
\protect\hypertarget{afghan-ux30c7ux30fcux30bfux306eux8aadux307fux8fbcux307f-1}{}
\begin{itemize}
\tightlist
\item
  アフガニスタンの調査データをRに読み込みます。
\end{itemize}

\scriptsize

\begin{Shaded}
\begin{Highlighting}[]
\CommentTok{\# 1. ローカルに保存したafghanデータの読み込み（推奨）}
\NormalTok{afghan }\OtherTok{\textless{}{-}} \FunctionTok{read.csv}\NormalTok{(}\StringTok{"afghan.csv"}\NormalTok{)}

\CommentTok{\# （参考）URLから直接読み込むことも可能}
\CommentTok{\# afghan \textless{}{-} read.csv("https://ayumu{-}tanaka.github.io/QSS/QSS\_Data/afghan.csv")}

\CommentTok{\# 2. データの次元（行数と列数）を確認}
\CommentTok{\# 2,754人の回答、10個の変数}
\FunctionTok{dim}\NormalTok{(afghan)}
\end{Highlighting}
\end{Shaded}

\begin{verbatim}
## [1] 2754   11
\end{verbatim}

\normalsize
\end{frame}

\begin{frame}[fragile]{回答者の基本属性の確認 (2)}
\protect\hypertarget{ux56deux7b54ux8005ux306eux57faux672cux5c5eux6027ux306eux78baux8a8d-2}{}
\begin{itemize}
\tightlist
\item
  \texttt{summary()}
  関数を使って、年齢や教育年数、収入などの分布を素早く把握します。
\end{itemize}

\scriptsize

\begin{Shaded}
\begin{Highlighting}[]
\CommentTok{\# 1. 年齢 (age) の要約}
\FunctionTok{summary}\NormalTok{(afghan}\SpecialCharTok{$}\NormalTok{age)}
\end{Highlighting}
\end{Shaded}

\begin{verbatim}
##    Min. 1st Qu.  Median    Mean 3rd Qu.    Max. 
##   15.00   22.00   30.00   32.39   40.00   80.00
\end{verbatim}

\begin{Shaded}
\begin{Highlighting}[]
\CommentTok{\# 2. 教育年数 (educ.years) の要約}
\FunctionTok{summary}\NormalTok{(afghan}\SpecialCharTok{$}\NormalTok{educ.years)}
\end{Highlighting}
\end{Shaded}

\begin{verbatim}
##    Min. 1st Qu.  Median    Mean 3rd Qu.    Max. 
##   0.000   0.000   1.000   4.002   8.000  18.000
\end{verbatim}

\begin{Shaded}
\begin{Highlighting}[]
\CommentTok{\# 3. 収入 (income) の要約}
\FunctionTok{summary}\NormalTok{(afghan}\SpecialCharTok{$}\NormalTok{income)}
\end{Highlighting}
\end{Shaded}

\begin{verbatim}
##    Length     Class      Mode 
##      2754 character character
\end{verbatim}

\normalsize
\end{frame}

\hypertarget{ux30abux30c6ux30b4ux30eaux5909ux6570ux306eux96c6ux8a08}{%
\section{3.1.2
カテゴリ変数の集計}\label{ux30abux30c6ux30b4ux30eaux5909ux6570ux306eux96c6ux8a08}}

\begin{frame}[fragile]{被害経験のクロス集計}
\protect\hypertarget{ux88abux5bb3ux7d4cux9a13ux306eux30afux30edux30b9ux96c6ux8a08}{}
\begin{itemize}
\tightlist
\item
  ISAF による被害経験とタリバンによる被害経験を同時に集計します。
\end{itemize}

\scriptsize

\begin{Shaded}
\begin{Highlighting}[]
\CommentTok{\# 1. table() で2変数のクロス集計を行う}
\CommentTok{\# 1 = 被害あり, 0 = 被害なし}
\NormalTok{victim.tab }\OtherTok{\textless{}{-}} \FunctionTok{table}\NormalTok{(}\AttributeTok{ISAF =}\NormalTok{ afghan}\SpecialCharTok{$}\NormalTok{violent.exp.ISAF, }
                    \AttributeTok{Taliban =}\NormalTok{ afghan}\SpecialCharTok{$}\NormalTok{violent.exp.taliban)}

\CommentTok{\# 2. prop.table() で割合（比率）に変換}
\FunctionTok{prop.table}\NormalTok{(victim.tab)}
\end{Highlighting}
\end{Shaded}

\begin{verbatim}
##     Taliban
## ISAF         0         1
##    0 0.4953445 0.1318436
##    1 0.1769088 0.1959032
\end{verbatim}

\normalsize

\begin{itemize}
\tightlist
\item
  \textbf{解釈}:
  両者から被害を受けた人、どちらか一方から受けた人、どちらからも受けていない人の割合がわかる。
\end{itemize}
\end{frame}

\hypertarget{ux307eux3068ux3081}{%
\section{3.1.3 まとめ}\label{ux307eux3068ux3081}}

\begin{frame}[fragile]{このセクションのまとめ}
\protect\hypertarget{ux3053ux306eux30bbux30afux30b7ux30e7ux30f3ux306eux307eux3068ux3081}{}
\begin{itemize}
\tightlist
\item
  \textbf{測定の重要性}:
  直接観察が困難な紛争地の被害実態を、大規模な調査データから数値化する。
\item
  \textbf{データの理解}:

  \begin{itemize}
  \tightlist
  \item
    数値変数の要約 (\texttt{summary()})
    により、中央値や平均、欠損値の有無を確認できる。
  \item
    カテゴリ変数の集計 (\texttt{table()}, \texttt{prop.table()})
    により、グループ間の関係を把握できる。
  \end{itemize}
\item
  \textbf{Rの操作}:

  \begin{itemize}
  \tightlist
  \item
    \texttt{read.csv()} による読み込み。
  \item
    \texttt{dim()} や \texttt{summary()} によるデータ構造の把握。
  \item
    \texttt{table()} と \texttt{prop.table()} を組み合わせた比率の算出。
  \end{itemize}
\end{itemize}
\end{frame}

\end{document}
